// String Format — demonstrate string functions and image statistics
//
// Educational example: compute basic image statistics (mean luminance,
// min/max channel values) across sampled points, then use format() to
// build a human-readable summary string.
//
// Output is a string, not an image. Use this with a string output node
// or for diagnostic/logging purposes.
//
// Connect @image to any RGB input.

s$label = "Image Stats";

// ── Sample image statistics from a grid ─────────────────────────────
// Sample a 5x5 grid to approximate global stats
float total_lum = 0.0;
float min_val = 1.0;
float max_val = 0.0;
float min_r = 1.0;
float max_r = 0.0;
float min_g = 1.0;
float max_g = 0.0;
float min_b = 1.0;
float max_b = 0.0;
int count = 0;

for (int sy = 0; sy < 5; sy++) {
    for (int sx = 0; sx < 5; sx++) {
        float su = (float(sx) + 0.5) / 5.0;
        float sv = (float(sy) + 0.5) / 5.0;
        vec3 s = sample(@image, su, sv);

        float lum = luma(s);
        total_lum += lum;

        min_val = min(min_val, min(s.r, min(s.g, s.b)));
        max_val = max(max_val, max(s.r, max(s.g, s.b)));
        min_r = min(min_r, s.r);
        max_r = max(max_r, s.r);
        min_g = min(min_g, s.g);
        max_g = max(max_g, s.g);
        min_b = min(min_b, s.b);
        max_b = max(max_b, s.b);
        count++;
    }
}

float mean_lum = total_lum / float(count);

// ── Build formatted output string ───────────────────────────────────
// format() fills {} placeholders in order (specs like {:.2f} work too) —
// it does not substitute % sequences.
string header = format("[{}]", $label);
string lum_info = format("Mean Luma: {}", mean_lum);
string range_info = format("Range: [{}, {}]", min_val, max_val);
string r_info = format("R: [{}, {}]", min_r, max_r);
string g_info = format("G: [{}, {}]", min_g, max_g);
string b_info = format("B: [{}, {}]", min_b, max_b);

// Build the report by concatenating lines
string nl = "\n";
string result = header + nl + lum_info + nl + range_info + nl + r_info + nl + g_info + nl + b_info;

@OUT = result;
